\section{Soundness from coupled comparisons}
\label{sec:soundness}

Assume that the game has value at most \(\theta\), and suppose that
the constructed digraph has a feedback vertex set \(F\) of cost at
most \(D\).  Continue with the fixed topological order \(\prec\)
and the conditioned functions from Sections~\ref{sec:ranks}
and~\ref{sec:prefix}.  In particular, \(\mathcal D\) denotes the
shared data in Equation~\eqref{eq:common-conditioning}, and a fresh
edge \(e=(u,v)\) has law \(\nu\) independently of \(\mathcal D\).
The random variables \(f_w,h_w(p)\), and \(B_w=B(f_w)\) are defined
for every vertex \(w\) once \(\mathcal D\) is fixed.

For a fixed \(i\in[M]\), there are two experiments for the fresh
subsets at the sampled edge.  In the \emph{coupled experiment},
\((R_u,R_v)\) has the law of the comparison test at index \(i\).
In the \emph{independent experiment}, conditionally on
\((\mathcal D,e)\), the subsets are independent with laws
\(\mu_{p_i}^{X}\) and \(\mu_{q_i}^{Y}\).
We write \(\Pr_{\mathrm c,i}\) and \(\Pr_{\mathrm{ind},i}\) for
probabilities in these experiments, including the common law of
\((\mathcal D,e)\).  Neither experiment samples a candidate cell
unless this is explicitly stated.

\subsection{Coupling and deletion cost}

Recall the coupling from the comparison family in
Definition~\ref{def:tests}.  For \(p=p_i\le1/4\), its exact
marginals are \(\mu_p^X\) and \(\mu_{2p}^Y\), and
Equation~\eqref{eq:coupling-inclusion} gives
\[
  R_u\subseteq\pi_e^{-1}(R_v)
  \quad\hbox{in every coupled outcome.}
\]
For each \(y\in Y\), the block consists of the two bits in
\(\pi_e^{-1}(y)\) and the bit at \(y\).  These blocks are
independent across \(y\).  The marginals and block independence
hold for every fixed edge, independently of the shared prefix data.

\begin{lemma}[Comparison deletion charge]
\label{lem:comparison-charge}
For every \(i\in[M]\),
\begin{equation}
\label{eq:coupled-comparison}
  \Pr_{\mathrm c,i}[\,f_u(R_u)=1,\ f_v(R_v)=0\,]
    \le 2\rho+\frac{M2^T D}{K}.
\end{equation}
\end{lemma}

\begin{proof}
Let \(\mathcal G\) be the event that both extended tuples
\(\mathbf s u\) and \(\mathbf s v\) are good.  The tuple choices
are independent of \(\psi\), so Lemma~\ref{lem:goodness}, averaged
over these choices, gives \(\Pr[\mathcal G^c]\le2\rho\).
On \(\mathcal G\), both conditioned functions use the same prefix
transition index \(r_*\), by Definition~\ref{def:conditioned-functions}
and Lemma~\ref{lem:compatibility}.

Suppose that \(\mathcal G\) holds and that
\(f_u(R_u)=1,\ f_v(R_v)=0\).  Add the candidate-cell draw from
the comparison test and consider the outcome \(r=r_*\).
The start and end rank vertices of Equation~\eqref{eq:comparison-path}
are then
\[
  b_{\mathbf s v}^{\psi}
      (E_v^{r_*}(R_v);2,17),
  \qquad
  b_{\mathbf s u}^{\psi}
      (E_u^{r_*}(R_u);3,18).
\]
The triples \((2,15,17)\) and \((3,15,18)\) are usable.
Goodness and the two function values show that the start does not
precede \(o_{\psi(15)}\), whereas the end precedes \(o_{\psi(15)}\).
If the test vertex were retained, its two arcs would require the
start to precede the test vertex and the test vertex to precede the
end.  This contradicts the fixed topological order.  Hence that
comparison vertex belongs to \(F\).

For fixed \((\mathcal D,e,R_u,R_v)\) as above, the candidate-cell
draw hits \(r_*\) with probability \(2^{-t}\ge2^{-T}\).
Let \(\delta_i\) be the probability that a comparison vertex is
deleted when its fresh threshold index is fixed to \(i\), including
the candidate-cell draw.  We have proved
\[
  \delta_i\ge
   2^{-T}\Pr_{\mathrm c,i}
     [\,\mathcal G,\ f_u(R_u)=1,\ f_v(R_v)=0\,].
\]
In the actual comparison family the threshold index is uniform on
\([M]\).  Its deleted weight is
\[
  \frac K M\sum_{j=1}^{M}\delta_j\le D,
\]
so \(\delta_i\le MD/K\).
Combining these inequalities and then adding
\(\Pr[\mathcal G^c]\le2\rho\) proves
Equation~\eqref{eq:coupled-comparison}.
\end{proof}

\subsection{Lists chosen before the edge}

\begin{lemma}[List decoding against game soundness]
\label{lem:list-decoding}
Suppose a game has value at most \(\theta\).
Assign to every game vertex \(w\) a list
\(J_w\subseteq\Lambda_w\), possibly empty, of cardinality at most
an integer \(J\ge1\).  Then
\[
  \Pr_{e=(u,v)\sim\nu}
    [\,\pi_e(J_u)\cap J_v\ne\varnothing\,]\le\theta J^2.
\]
The same conclusion holds conditionally on any data independent of
the fresh edge, if all the lists are fixed by those data.
\end{lemma}

\begin{proof}
Independently for each vertex with a nonempty list, choose its label
uniformly from that list.  Give every empty-list vertex any fixed
label from its alphabet.  If
\(\pi_e(J_u)\cap J_v\ne\varnothing\), there is at least one pair
\((x,y)\in J_u\times J_v\) with \(\pi_e(x)=y\).  Since \(u\) and
\(v\) are different vertices, the probability of choosing that pair
is
\[
  \frac1{\abs{J_u}\abs{J_v}}\ge\frac1{J^2}.
\]
The expected value of this random labeling is therefore at least
the displayed intersection probability divided by \(J^2\).
Every deterministic labeling in its support has value at most
\(\theta\), so the expected value is at most \(\theta\).
This proves the claim.  For the conditional assertion, fix the
conditioning data first and repeat the same argument; independence
leaves the fresh edge law equal to \(\nu\).
\end{proof}

\subsection{Replacing the coupling by product measure}

\begin{proposition}[Product comparison]
\label{prop:product-comparison}
For every \(i\in[M]\), under the common distribution of
\((\mathcal D,e)\),
\begin{equation}
\label{eq:product-comparison}
  \E\bigl[h_u(p_i)(1-h_v(q_i))\bigr]\le\frac{p_i}{16}.
\end{equation}
\end{proposition}

\begin{proof}
Fix \(i\).  For each value of \(\mathcal D\), construct a junta
approximant for each game vertex, before drawing the fresh edge.
For notation, put
\[
  \beta_{w,i}=
  \begin{cases}
    p_i,&w\in U,\\
    q_i,&w\in V.
  \end{cases}
\]
If \(B_w\le B_0\), Lemma~\ref{lem:pivotal-inequalities} gives
\(I_{\beta_{w,i}}(f_w)\le B_0\).
Corollary~\ref{cor:dyadic-junta} and the choice of \(J\) in
Section~\ref{sec:parameters} therefore provide a Boolean function
\(\widetilde f_{w,i}\) on a coordinate set
\(J_w\subseteq\Lambda_w\), with
\begin{equation}
\label{eq:junta-error}
  \abs{J_w}\le J,\qquad
  \Pr_{R\sim\mu_{\beta_{w,i}}^{\Lambda_w}}
    [\,f_w(R)\ne\widetilde f_{w,i}(R)\,]\le\gamma.
\end{equation}
Choose one such pair by a fixed ordering of the finitely many
possible coordinate sets and Boolean truth tables.
If \(B_w>B_0\), set \(J_w=\varnothing\) and choose
\(\widetilde f_{w,i}\) to be the constant-zero function; no error
guarantee is asserted at that vertex.
We suppress the dependence of \(J_w\) on \(\mathcal D\) and \(i\).

This defines one list for every actual game vertex from
\((\mathcal D,i)\) alone.  The functions \(f_w\), including the
fallback functions, are already defined for all vertices from
\(\mathcal D\) alone.  Thus the choice of a junta or list uses no
information about the fresh sampled edge.  It may depend on the
fixed game instance and the fixed topological order, as may any
labeling used in a soundness argument.  A repetition in the prefix,
or an endpoint that also occurs in the prefix, does not change this
conditional independence: the fresh edge was drawn independently
of the entire prefix.

By Lemma~\ref{lem:list-decoding}, for every fixed value of
\(\mathcal D\),
\[
  \Pr_{e\sim\nu}
     [\,\pi_e(J_u)\cap J_v\ne\varnothing
        \mid\mathcal D\,]\le\theta J^2.
\]
After averaging over \(\mathcal D\), this remains the bound on the
list-intersection event under the common law.
Also, Proposition~\ref{prop:pivotal-budget} and Markov's
inequality give the averaged bound
\begin{equation}
\label{eq:budget-cutoff}
  \Pr[\,B_u>B_0\ \text{or}\ B_v>B_0\,]
   \le\frac{\E B_u+\E B_v}{B_0}
   \le 2\gamma.
\end{equation}
Here we have used the actual endpoint marginals of \(\nu\).
Equation~\eqref{eq:budget-cutoff} is a bound under the joint law of
\((\mathcal D,e)\); no corresponding pointwise bound conditional
on \(\mathcal D\) is needed.

Let \(\mathcal A_i\) be the event, determined by
\((\mathcal D,e)\), that both budgets are at most \(B_0\) and
\(\pi_e(J_u)\cap J_v=\varnothing\).  The preceding bounds imply
\begin{equation}
\label{eq:junta-exceptions}
  \Pr[\mathcal A_i^c]\le2\gamma+\theta J^2.
\end{equation}
Fix \((\mathcal D,e)\) in \(\mathcal A_i\).
In the coupled experiment, all bits used by the big-side junta
belong to fibers indexed by \(\pi_e(J_u)\), and all bits used by
the small-side junta belong to fibers indexed by \(J_v\).
These are disjoint collections of independently generated fibers.
Consequently
\(\widetilde f_{u,i}(R_u)\) and
\(\widetilde f_{v,i}(R_v)\) are independent in the coupled
experiment.  They are also independent in the independent
experiment, and their marginal distributions are unchanged.
Their joint distributions therefore agree in the two experiments.
This includes the case of a constant junta with an empty
coordinate set.

The conditional approximation error for each endpoint in
Equation~\eqref{eq:junta-error} is at most \(\gamma\) in either
experiment, because both have exactly the stipulated product
marginals.  Replacing \(f_u,f_v\) by their respective juntas changes
the conditional probability of the pattern \((1,0)\) by at most
\(2\gamma\), by a union bound.  Doing this in both experiments and
using equality of the junta joint distributions yields
\begin{align*}
 &\left|
  \Pr_{\mathrm{ind},i}[\,f_u(R_u)=1,\ f_v(R_v)=0
                      \mid\mathcal D,e\,]
 \right.\\[-2mm]
 &\hspace{22mm}\left.
  -\Pr_{\mathrm c,i}[\,f_u(R_u)=1,\ f_v(R_v)=0
                      \mid\mathcal D,e\,]
  \right|\le4\gamma
\end{align*}
on \(\mathcal A_i\).
On its complement the difference is at most one.
Integrating over the unchanged law of \((\mathcal D,e)\), and using
Equation~\eqref{eq:junta-exceptions} and
Lemma~\ref{lem:comparison-charge}, gives
\begin{align}
\label{eq:product-comparison-error}
  &\Pr_{\mathrm{ind},i}[\,f_u(R_u)=1,\ f_v(R_v)=0\,]\notag\\
  &\qquad\le
  2\rho+\frac{M2^T D}{K}
      +2\gamma+\theta J^2+4\gamma
  \le10\gamma.
\end{align}
The conditional probability of the pattern in the independent
experiment is exactly \(h_u(p_i)(1-h_v(q_i))\).
Finally, \(\gamma=p_M/1000\) and \(p_M\le p_i\), so
\[
  10\gamma=\frac{p_M}{100}\le\frac{p_i}{16}.
\]
This proves Equation~\eqref{eq:product-comparison}.
\end{proof}

\subsection{The dyadic tail contradiction}

\begin{proposition}[Weighted soundness]
\label{prop:weighted-soundness}
If the game has value at most \(\theta\), every feedback vertex set
of the constructed weighted digraph has cost strictly greater
than \(D\).
\end{proposition}

\begin{proof}
Suppose that \(F\) has cost at most \(D\), as assumed throughout
this section.  Fix \(i\in[M]\), and let
\[
  \mathcal H_i=\{h_v(q_i)\ge1/2\}
\]
under the common distribution of \((\mathcal D,e)\).
On its complement,
\(h_u(p_i)\le2h_u(p_i)(1-h_v(q_i))\).  Consequently
Propositions~\ref{prop:prefix-lower-bound}
and~\ref{prop:product-comparison} imply
\begin{align*}
  \E[\,h_u(p_i)\ind{\mathcal H_i}\,]
  &=\E h_u(p_i)
       -\E[\,h_u(p_i)\ind{\mathcal H_i^c}\,]\\
  &\ge\frac{p_i}{4}
       -2\E[\,h_u(p_i)(1-h_v(q_i))\,]\\
  &\ge\frac{p_i}{8}.
\end{align*}
Since \(0\le h_u(p_i)\le1\), we obtain
\(\Pr[\mathcal H_i]\ge p_i/8\).
Lemma~\ref{lem:pivotal-inequalities} also gives
\(h_v(q_i)\le q_iB_v=2p_iB_v\) pointwise.
Thus, on \(\mathcal H_i\), one has \(B_v\ge1/(4p_i)\), and hence
\begin{equation}
\label{eq:dyadic-budget-tails}
  \Pr\!\left[B_v\ge\frac1{4p_i}\right]\ge\frac{p_i}{8},
  \qquad i\in[M].
\end{equation}

These are tail estimates for one random variable \(B_v\) under one
fixed probability law.  The definition of \(f_v\), and therefore
that of \(B_v\), does not use the fresh threshold index \(i\).
The junta approximants used to prove the individual estimates may
depend on \(i\), but this does not alter \(B_v\).

Put \(a_i=1/(4p_i)\) for \(i\ge1\), and \(a_0=0\).
For every real \(x\ge0\),
\[
  x\ge\sum_{i=1}^{M}(a_i-a_{i-1})\ind{x\ge a_i}.
\]
Indeed, if \(a_j\) is the largest crossed threshold, the sum on the
right telescopes to \(a_j\le x\); if no threshold is crossed, it
is zero.  Taking expectations and using
Equation~\eqref{eq:dyadic-budget-tails} yields
\[
  \E B_v
  \ge\sum_{i=1}^{M}(a_i-a_{i-1})\Pr[B_v\ge a_i]
  \ge\sum_{i=1}^{M}\frac{a_i}{2}\frac{p_i}{8}
  =\frac M{64}.
\]
Here \(a_i=2a_{i-1}\) for \(i\ge2\), while
\(a_1-a_0=a_1\ge a_1/2\).
Since \(M=128(D+1)\), this gives
\(\E B_v\ge2(D+1)\), contradicting
Proposition~\ref{prop:pivotal-budget}.
No feedback vertex set of cost at most \(D\) can therefore exist.
\end{proof}
